\subsection{Conclusion: Toward the AI Scientist \quad / \quad \textthai{บทสรุป: มุ่งสู่นักวิทยาศาสตร์ AI}}
\begin{paracol}{2}
We have traveled from the simple oscillations of a pendulum to the complex collisions of subatomic particles. Throughout this book, we have seen that AI is more than just a data-fitting tool; it is a way to encode physical intuition into mathematical structures.

In the next decade, we will see the rise of the \textbf{Autonomous Scientific Collaborator}\index{Collaborator@ผู้ร่วมงาน (Collaborator)}—systems that can hypothesize, design an experiment, and execute it via robotics in a closed loop. The goal of this book was to give you the "Keys to the Lab" for that future.

\switchcolumn

\begin{thai}
เราได้เดินทางตั้งแต่การแกว่งที่เรียบง่ายของลูกตุ้มไปจนถึงการชนที่ซับซ้อนของอนุภาคระดับต่ำกว่าอะตอม ตลอดทั้งเล่มนี้ เราได้เห็นแล้วว่า AI เป็นมากกว่าเครื่องมือจับคู่ข้อมูล แต่มันคือวิธีในการเข้ารหัสสัญชาตญาณทางฟิสิกส์ให้กลายเป็นโครงสร้างทางคณิตศาสตร์

ในทศวรรษหน้า เราจะเห็นการเกิดขึ้นของ \textbf{ผู้ร่วมงานทางวิทยาศาสตร์อัตโนมัติ} - ระบบที่สามารถตั้งสมมติฐาน ออกแบบการทดลอง และดำเนินการผ่านหุ่นยนต์ในลูปปิด เป้าหมายของหนังสือเล่มนี้คือการมอบ "กุญแจสู่ห้องแล็บ" สำหรับอนาคตนั้นให้แก่คุณ
\end{thai}
\end{paracol}

\subsection{The Ethics of Discovery \quad / \quad \textthai{จริยธรรมแห่งการค้นพบ}}
\begin{paracol}{2}
As we reach this frontier, we must address the **Ethical Responsibilities**\index{Ethics@จริยธรรม (Ethics)} of the physicist. We must strive for \textbf{Interpretability}—refusing to accept an AI discovery unless we can extract the underlying physical law. We must also ensure \textbf{Open Science}, where the tools of discovery are accessible to all, preventing a "Computational Divide" in human knowledge.

\switchcolumn

\begin{thai}
เมื่อเรามาถึงพรมแดนนี้ เราต้องจัดการกับ **ความรับผิดชอบทางจริยธรรม** ของมนุษยชาติ เราต้องมุ่งมั่นเพื่อ \textbf{ความสามารถในการตีความ}—โดยปฏิเสธการยอมรับการค้นพบของ AI เว้นแต่เราจะสามารถดึงกฎทางฟิสิกส์พื้นฐานออกมาได้ เรายังต้องรับประกัน \textbf{วิทยาศาสตร์แบบเปิด} ซึ่งเครื่องมือในการค้นหาสามารถเข้าถึงได้โดยทุกคน เพื่อป้องกัน "ช่องว่างทางคอมพิวเตอร์" ในความรู้ของมนุษย์
\end{thai}
\end{paracol}

\begin{aibox}{Final Thought: A Call to Discovery}
"The best way to predict the future of physics is to program it." — The stars are no longer the limit; they are the next dataset.
\end{aibox}
